\documentclass[10pt]{article}

\usepackage[margin=0.75in]{geometry}
\usepackage{amsmath,amsthm,amssymb}
\usepackage{xcolor}
\usepackage{cancel}
\usepackage{graphicx}
\usepackage{changepage}
\usepackage{circuitikz}
\usepackage{pgfplots}
\usepackage{physics}
\usepackage{hyperref}
\usepackage{siunitx}
\usepackage{minted}
\usepackage{fontspec}
\usepackage{relsize}
\usepackage{subfig}
\usepackage{mhchem}
\usepackage{todonotes}
\usepackage{biblatex}
\usepackage{multicol, multirow, booktabs}
\usepackage[breakable]{tcolorbox}
\usepackage[inline]{enumitem}
\patchcmd{\thebibliography}{\section*{\refname}}{}{}{}
\addbibresource{sources.bib}

\theoremstyle{definition}
\newtheorem{problem}{Problem}
\newtheorem{soln}{Solution}

\pgfplotsset{compat=newest}
\usetikzlibrary{lindenmayersystems}
\usetikzlibrary{arrows}
\usetikzlibrary{calc}
\usetikzlibrary{positioning, fit}
\usetikzlibrary{3d, perspective}

\definecolor{incolor}{HTML}{303F9F}
\definecolor{outcolor}{HTML}{D84315}
\definecolor{cellborder}{HTML}{CFCFCF}
\definecolor{cellbackground}{HTML}{F7F7F7}
\newcommand{\ui}{\hat{i}}
\newcommand{\uj}{\hat{j}}
\newcommand{\uk}{\hat{k}}
\newcommand{\ux}{\hat{x}}
\newcommand{\uy}{\hat{y}}
\newcommand{\uz}{\hat{z}}
\newcommand{\uv}[1]{\hat{\mathbf{{#1}}}}
\newcommand{\pr}[1]{{#1^\prime}}
\newcommand{\justif}[2]{&{#1}&\text{#2}}
\newcommand{\dul}[1]{\underline{\underline{#1}}}
\pgfdeclarelayer{background}  
\pgfsetlayers{background,main}
\AtBeginDocument{\RenewCommandCopy\qty\SI}

\makeatletter
\newcommand{\boxspacing}{\kern\kvtcb@left@rule\kern\kvtcb@boxsep}
\makeatother
\newcommand{\prompt}[4]{
    \ttfamily\llap{{\color{#2}[#3]:\hspace{3pt}#4}}\vspace{-\baselineskip}
}

\newcommand{\thevenin}[2]{
  \begin{center}
    \begin{circuitikz} \draw
      (0,0) -- (2,0) to[battery1, l_=$V_{Th}\eq#1$] (2,2) 
      to[resistor, l_=$R_{Th}\eq#2$] (0,2)
      ;
      \draw [o-] (-.07,2.079);
      \draw [o-] (-.07,0.079);
    \end{circuitikz}
  \end{center}
}

\newcommand{\norton}[2]{
  \begin{center}
    \begin{circuitikz} \draw
      (0,0) -- (3,0) to[american current source, l_=$I_{N}\eq#1$] (3,2) -- (0,2) (2,0)
      to[resistor, l=$R_{N}\eq#2$] (2,2)
      ;
      \draw [o-] (-.07,2.079);
      \draw [o-] (-.07,0.079);
    \end{circuitikz}
  \end{center}
}

\newcommand{\highlight}[1]{\colorbox{yellow}{$\displaystyle #1$}}

\newcommand{\ti}[1]{\widetilde{#1}}

\newfontface{\Kaufmann}{Kaufmann}
\DeclareTextFontCommand{\kf}{\Kaufmann}
\newcommand{\scriptr}{\fontsize{12pt}{12pt}\kf{r}}

\newfontface{\KaufmannB}{Kaufmann Bold}
\DeclareTextFontCommand{\kfb}{\KaufmannB}
\newcommand{\bscriptr}{\fontsize{12pt}{12pt}\kfb{r}}

\newcommand{\bv}[1]{\vec{#1}}

\title{Physics 4605H: Assignment II}
\author{Jeremy Favro (0805980) \\ Trent University, Peterborough, ON, Canada}
\date{\today}

\begin{document}
\maketitle

% PROBLEM 1
\begin{problem}
Find an article which discusses quantum computing in a newspaper or magazine (including 
online equivalents). What does the article say about possible applications for quantum
computing? Post on the Blackboard discussion forum:
\begin{enumerate}[label=(\roman*)]
  \item a reference and/or a link to the
        article;
  \item your rough assessment of the general reliability of the source;
  \item at least one
        possible application of quantum computing mentioned in the article.
\end{enumerate}
\end{problem}
\begin{soln}
  Done!
\end{soln}

% PROBLEM 2 
\begin{problem}
Use the CRC Handbook of Chemistry and Physics (available online through Bata) to
look up the following information:
\begin{enumerate}[label=(\alph*)]
  \item What is a nucleus (other than \ce{^{1}H} and \ce{^{13}C}) that has spin $1/2$?
  \item What is the gyromagnetic ratio of \ce{^{31}P}?
\end{enumerate}
In both cases, include where in the handbook you located the information you've used as
well as any calculations you've done.
\end{problem}
\begin{soln}
  \begin{enumerate}[label=(\alph*)]
    \item \ce{^{3}He}. I know it appears on the first page but I swear I'm not lazy. I was reading about \ce{^{3}He} fusion recently and so it was in my head. \cite{CRC}
    \item The CRC provides only the nuclear magnetic moment, defined to be ``Magnetic dipoles, in units of nuclear magneton (nm)''. The value of this associated with
          \ce{^{31}P} is 1.130925 \cite{CRC}. The nuclear magneton, which this quantity is in units of, is defined as $$\frac{e\hbar}{2m_p}$$
          where $e$ is the electron charge and $m_p$ the proton mass, both in SI units. Recalling that the gyromagnetic ratio $\gamma$ is present in the expression
          $$\mu_z=\gamma S_z$$
          and we have $\mu_z/\mu_N=1.130925\implies 1.13\mu_N=\mu_z$, we get
          $$\gamma \approx 2\cdot1.13\mu_N/h$$
          where the division by $h$ is to obtain the final answer in the desired units of $\unit{\mega\hertz\per\tesla}$.
          I used a program called Qalculate to evaluate this which handles the units for me and gives a result of
          $\gamma=\qty{17.2411625}{\mega\hertz\per\tesla}$
  \end{enumerate}
\end{soln}

% PROBLEM 3
\begin{problem}
Consider a spin-$1/2$ object in a uniform magnetic field $\bv{B}=B_0\uz$. Ignore the chemical
shift, such that
$$\hat{H}=-\bv{\mu}\cdot\bv{B}=-\frac{\gamma \hbar B_0}{2}\hat{\sigma}_z.$$
Let the initial state of the particle be
$$\ket{\psi(t=0)}=\frac{1}{\sqrt{2}}\ket{0}+\frac{1}{\sqrt{2}}\ket{1}$$
\begin{enumerate}[label=(\alph*)]
  \item Show that this state is normalized.
  \item Show that this is an eigenstate of the operator $\hat{S}_x$ corresponding to eigenvalue $+\hbar/2$.
  \item Show that the expectation value of $S_x$ in this state is $+\hbar/2$.
  \item What is the expectation value of $S_y$ in this state?
  \item What is the expectation value of $S_z$ in this state?
\end{enumerate}

\noindent Now consider the time dependence. Recall that $\ket{0}$ and $\ket{1}$ are eigenstates of the Hamiltonian.
\begin{enumerate}[label=(\alph*)]
  \setcounter{enumi}{5}
  \item What is the wavefunction at a later time $t$?
  \item What is the expectation value of $S_x$ as a function of time?
  \item What is the expectation value of $S_y$ as a function of time?
  \item What is the expectation value of $S_z$ as a function of time?
  \item At time $t = 0$, as shown in (c)-(e), the expectation value of $\bv{S}$ is entirely along $x$. What
        is the (first) time at which it will be entirely along $y$?
  \item In what direction does it rotate?
\end{enumerate}

\end{problem}
\begin{soln}
  \begin{enumerate}[label=(\alph*)]
    \item \begin{align*}
            \bra{\psi(t=0)}\ket{\psi(t=0)} & =\left[\frac{1}{\sqrt{2}}\bra{0}+\frac{1}{\sqrt{2}}\bra{1}\right]\left[\frac{1}{\sqrt{2}}\ket{0}+\frac{1}{\sqrt{2}}\ket{1}\right] \\
                                           & =\frac{1}{2}\bra{0}\ket{0}+\frac{1}{2}\bra{1}\ket{1}+\frac{1}{2}\bra{0}\ket{1}+\frac{1}{2}\bra{1}\ket{0}                          \\
                                           & =\frac{1}{2}+\frac{1}{2}=1 \justif{\quad}{Assuming the basis is orthonormal}
          \end{align*}
    \item We can write this state as $$\psi=\frac{1}{\sqrt{2}}\left[\begin{pmatrix}
                1 \\0
              \end{pmatrix}+\begin{pmatrix}
                0 \\1
              \end{pmatrix}\right]=\frac{1}{\sqrt{2}}\begin{pmatrix}
              1 \\1
            \end{pmatrix}.$$
          Recalling that our operator $\hat{S}_x$ is, in matrix representation, $$\dul{S}_x=\frac{\hbar}{2}\begin{pmatrix}
              0 & 1 \\
              1 & 0
            \end{pmatrix}$$
            we can evaluate the eigenvalue equation as
            $$\frac{\hbar}{2}\begin{pmatrix}
              0 & 1 \\
              1 & 0
            \end{pmatrix}\frac{1}{\sqrt{2}}\begin{pmatrix}
              1 \\1
            \end{pmatrix}=\frac{\hbar}{2}\frac{1}{\sqrt{2}}\begin{pmatrix}
              1 \\1
            \end{pmatrix}$$
            hence this is an eigenstate of $\hat{S}_x$ with eigenvalue $\hbar/2$.
    \item $$\langle S_x\rangle=\bra{\psi}\hat{S}_x\ket{\psi}=\frac{1}{\sqrt{2}}\begin{pmatrix}
              1 &1
            \end{pmatrix}\frac{\hbar}{2}\frac{1}{\sqrt{2}}\begin{pmatrix}
              1 \\1
            \end{pmatrix}=\frac{\hbar}{4}\cdot (1\cdot 1 +1\cdot 1)=\hbar/2
            $$
    \item Recall that we defined $$\dul{S}_y=\frac{\hbar}{2}\begin{pmatrix}
      0 & -i\\
      i & 0
    \end{pmatrix}$$
    and so $$\hat{S}_y\ket{\psi}=\frac{\hbar}{2}\begin{pmatrix}
      0 & -i\\
      i & 0
    \end{pmatrix}\frac{1}{\sqrt{2}}\begin{pmatrix}
              1 \\1
            \end{pmatrix}=\frac{i\hbar}{2\sqrt{2}}\begin{pmatrix}
              -1\\1
            \end{pmatrix}$$
            and so 
            $$\langle S_y\rangle=\bra{\psi}\hat{S}_y\ket{\psi}=\frac{1}{\sqrt{2}}\begin{pmatrix}
              1 &1
            \end{pmatrix}\frac{i\hbar}{2\sqrt{2}}\begin{pmatrix}
              -1\\1
            \end{pmatrix}=0$$
    \item Recall we defined 
            $$\dul{S}_z=\frac{\hbar}{2}\begin{pmatrix}
              1 & 0 \\
              0 & -1
            \end{pmatrix}$$
            so
            $$\hat{S}_z\ket{\psi}=\frac{\hbar}{2}\begin{pmatrix}
              1 & 0 \\
              0 & -1
            \end{pmatrix}\frac{1}{\sqrt{2}}\begin{pmatrix}
              1 \\1
            \end{pmatrix}=\frac{\hbar}{2\sqrt{2}}\begin{pmatrix}
              1 \\-1
            \end{pmatrix}$$
            and so 
            $$\langle S_z\rangle=\frac{1}{\sqrt{2}}\begin{pmatrix}
              1 &1
            \end{pmatrix}\frac{\hbar}{2\sqrt{2}}\begin{pmatrix}
              1 \\-1
            \end{pmatrix}=0$$
    \item What is the wavefunction at a later time $t$?
    \item What is the expectation value of $S_x$ as a function of time?
    \item What is the expectation value of $S_y$ as a function of time?
    \item What is the expectation value of $S_z$ as a function of time?
    \item At time $t = 0$, as shown in (c)-(e), the expectation value of $\bv{S}$ is entirely along $x$. What
          is the (first) time at which it will be entirely along $y$?
    \item In what direction does it rotate?
  \end{enumerate}
\end{soln}

\section{References}
\printbibliography
\end{document}